%! Radicals and Rational Exponents — Unit 1, Lesson 1.2 — Slides (Beamer)
\documentclass[aspectratio=169, 11pt]{beamer}
\usepackage{atda-beamer}   % provides \royalheader{} and \sectionlabel[color]{}

\begin{document}

% TITLE SLIDE (CourseName is not defined in beamer, so the name is literal)
{
\setbeamercolor{background canvas}{bg=royal}
\begin{frame}[plain]
  \vspace{2.8cm}\hspace{0.55cm}%
  \begin{minipage}{0.9\textwidth}
    {\color{goldacc}\bfseries\small LESSON 1.2}\\[0.5cm]
    {\color{white}\bfseries\LARGE Radicals and Rational Exponents}\\[1.2cm]
    {\color{mistmid}\small Algebra, Trigonometry, and Data Analysis~$\cdot$~Unit 1}
  \end{minipage}
\end{frame}
}

% ── HOOK ──────────────────────────────────────────────────────────────────────
\begin{frame}[t, plain]
  \royalheader{Reading a Skid Mark}
  \sectionlabel[goldacc]{hook --- vote before you compute}
  Crash investigators estimate speed from skid length. On a road with drag factor $f$, a skid of
  $d$ feet means the car was going about
  \[ s = \sqrt{30fd} \quad \text{mph.} \]
  Same dry asphalt ($f=0.8$). Car A skids $60$ ft. Car B skids $240$ ft --- \textbf{four times}
  as far.
  \begin{block}{Was Car B going four times as fast?}
    Vote first. Then we compute.
  \end{block}
\end{frame}

% ── HOOK PAYOFF ───────────────────────────────────────────────────────────────
\begin{frame}[t, plain]
  \royalheader{Twice, Not Four Times}
  \sectionlabel{a square root undoes what an exponent did}
  \[
  \begin{aligned}
    \text{Car A:} \quad \sqrt{30(0.8)(60)}  &= \sqrt{1440} = 12\sqrt{10} \approx 37.9 \\
    \text{Car B:} \quad \sqrt{30(0.8)(240)} &= \sqrt{5760} = 24\sqrt{10} \approx 75.9
  \end{aligned}
  \]
  \begin{block}{Four times the skid, twice the speed}
    Because $\sqrt{4}=2$. Last lesson doubling a radius \emph{quadrupled} the area; this lesson
    runs that machinery backwards.
  \end{block}
\end{frame}

% ── SIMPLIFYING ───────────────────────────────────────────────────────────────
\begin{frame}[t, plain]
  \royalheader{Simplifying: Pull Out the Perfect Powers}
  \sectionlabel{find the largest one --- then there is nothing left to pull}
  \[ \sqrt[n]{ab}=\sqrt[n]{a}\cdot\sqrt[n]{b} \qquad
     \sqrt[n]{\tfrac{a}{b}}=\tfrac{\sqrt[n]{a}}{\sqrt[n]{b}} \]
  \renewcommand{\arraystretch}{1.5}
  \begin{tabularx}{\textwidth}{@{}p{3.4cm} p{4.6cm} X@{}}
    \textbf{Radical} & \textbf{Split it} & \textbf{Simplified} \\
    \hline
    $\sqrt{50}$       & $\sqrt{25\cdot 2}$       & $5\sqrt{2}$ \\
    $\sqrt[3]{54}$    & $\sqrt[3]{27\cdot 2}$    & $3\sqrt[3]{2}$ \\
    $\sqrt{72x^5}$    & $\sqrt{36x^4\cdot 2x}$   & $6x^2\sqrt{2x}$ \\
  \end{tabularx}
  \vspace{0.3cm}

  {\color{goldacc}\bfseries $\sqrt{72}=\sqrt4\cdot\sqrt{18}$ is true --- but not simplified.}
\end{frame}

% ── RATIONAL EXPONENTS ────────────────────────────────────────────────────────
\begin{frame}[t, plain]
  \royalheader{Fractional Exponents Are Forced Too}
  \sectionlabel{the product rule from 1.1 decides this, not us}
  \[ 9^{1/2}\cdot 9^{1/2} = 9^{\frac12+\frac12} = 9^1 = 9 \]
  So $9^{1/2}$ must be the number that squares to $9$. That is what $\sqrt{9}$ means.
  \begin{itemize}
    \setlength{\itemsep}{0.3cm}
    \item $a^{1/n}=\sqrt[n]{a}$ --- the \textbf{denominator} is the index.
    \item $a^{m/n}=\sqrt[n]{a^m}=\left(\sqrt[n]{a}\right)^m$ --- the \textbf{numerator} is the
          power. Root first; the numbers stay smaller.
  \end{itemize}
  \vspace{0.2cm}

  {\color{royal}\bfseries $27^{2/3}=\left(\sqrt[3]{27}\right)^2=9$. \quad
   $16^{-1/2}=\tfrac14$ --- still not negative.}
\end{frame}

% ── EVERY 1.1 RULE SURVIVES ───────────────────────────────────────────────────
\begin{frame}[t, plain]
  \royalheader{Every Rule from Lesson 1.1 Still Holds}
  \sectionlabel{the exponents just became fractions}
  \[
  \begin{aligned}
    \left(32x^{10}\right)^{3/5} &= 32^{3/5}\left(x^{10}\right)^{3/5} \\
                                &= \left(\sqrt[5]{32}\right)^3 x^{6} \\
                                &= 8x^6
  \end{aligned}
  \]
  \begin{block}{Nothing new is happening here}
    Power of a product, then power of a power --- the same two rules you filled in yesterday.
    Only the exponent looks different.
  \end{block}
\end{frame}

% ── OPERATIONS ────────────────────────────────────────────────────────────────
\begin{frame}[t, plain]
  \royalheader{Adding Radicals Is the Like-Terms Rule}
  \sectionlabel{$\sqrt3$ behaves exactly like $x$}
  \[
  \begin{aligned}
    3\sqrt{12}+\sqrt{75} &= 3\left(2\sqrt3\right)+5\sqrt3 \;=\; 6\sqrt3+5\sqrt3 \;=\; 11\sqrt3 \\[4pt]
    \left(2\sqrt5\right)\left(3\sqrt{10}\right) &= 6\sqrt{50} \;=\; 30\sqrt2 \\[4pt]
    \left(\sqrt7+3\right)\left(\sqrt7-2\right)  &= 7-2\sqrt7+3\sqrt7-6 \;=\; 1+\sqrt7
  \end{aligned}
  \]
  \begin{block}{Simplify \emph{before} you decide they are unlike}
    $\sqrt{12}$ and $\sqrt{75}$ look unrelated. They are both $\sqrt3$ in disguise.
  \end{block}
\end{frame}

% ── RATIONALIZING ─────────────────────────────────────────────────────────────
\begin{frame}[t, plain]
  \royalheader{No Radical Downstairs}
  \sectionlabel{multiply by a form of 1 --- the value never moves}
  \[
  \begin{aligned}
    \dfrac{6}{\sqrt3} &= \dfrac{6}{\sqrt3}\cdot\dfrac{\sqrt3}{\sqrt3}
                       = \dfrac{6\sqrt3}{3} = 2\sqrt3
  \end{aligned}
  \]
  \begin{block}{The two errors to name out loud now}
    $\sqrt{9+16}$ is $5$, not $7$ --- a root splits over a \textbf{product}, never a sum.\\[2pt]
    And $16^{1/2}$ is $4$, not $8$ --- the exponent is a root, not ``divide by two.''
  \end{block}
\end{frame}

% ── ACTIVITY LAUNCH ───────────────────────────────────────────────────────────
\begin{frame}[t, plain]
  \royalheader{Group Activity: The Drop Tower}
  \sectionlabel{groups of three --- everyone starts at Tier R}
  An object dropped from rest falls $h=16t^2$ feet in $t$ seconds.
  \begin{block}{Your job}
    Solve $h=16t^2$ for $t$ to get $t=\dfrac{\sqrt h}{4}$. Find the fall time for a $128$-foot
    tower --- then tell marketing how tall the tower must be to \textbf{double} it.
  \end{block}
  \vspace{0.2cm}

  {\color{goldacc}\bfseries Predict the answer before you compute it. Most people guess wrong.}
\end{frame}

% ── CLOSE ─────────────────────────────────────────────────────────────────────
\begin{frame}[t, plain]
  \royalheader{Exit Ticket}
  \sectionlabel[goldacc]{notes closed --- 5 minutes}
  \begin{enumerate}
    \setlength{\itemsep}{0.35cm}
    \item Simplify $\sqrt{98x^7}$.
    \item Combine and simplify $\sqrt{27}+\sqrt{12}$.
    \item Write $8^{2/3}$ in radical form and evaluate it. What does the $3$ do? The $2$?
  \end{enumerate}
  \vspace{0.3cm}
  {\color{royal}\bfseries Next: Lesson 1.3 --- Factoring: GCF and Grouping.}
\end{frame}

\end{document}
